\section{Introduction}
\label{sec:introduction}

Species recognition is an unusually demanding image-classification problem.
The label space is large, distinctions may depend on a small anatomical
marking, and photographs are acquired opportunistically rather than under
controlled laboratory conditions.  In community-science collections, a
target can occupy only a few pixels, be partially occluded, or resemble a
closely related species while the background offers a tempting but unstable
cue.  The iNaturalist benchmarks capture these challenges at scale and have
therefore become important tests of visual representation learning
\cite{vanhorn2018inaturalist,vanhorn2021benchmarking}.

This work asks how far progressively more expressive computer-vision methods
can address a fixed 500-class recognition task.  We begin with interpretable
handcrafted representations---colour histograms, local binary patterns (LBP),
histograms of oriented gradients (HOG), and SIFT bag-of-visual-words
(BoVW)---paired with linear or tree-based classifiers.  We then compare
ResNet18 trained from random initialization against ImageNet-pretrained
ResNet18 and EfficientNet-B0 models.  Beyond clean-set accuracy, we test
controlled image corruptions and inspect decisions with Grad-CAM.  This broad
comparison is useful because a single accuracy number cannot reveal whether a
model is learning transferable morphology, memorising texture, or failing
under a realistic loss of image quality.

Equally important is evaluation discipline.  The official iNaturalist
validation archive is used only as a held-out test source.  Every fitted
transform, visual vocabulary, classifier, and neural-network update sees the
training partition alone.  Hyperparameters, early stopping, and ablations use
an internal validation partition.  Stage~3 records its selected configurations
before test access.  For the rewritten Stage~2 track, the present selection
manifest cryptographically freezes the complete selected state going forward,
but the available history cannot independently establish that an equivalent
lock existed before the already-recorded test evaluation.  We therefore do not
claim a retrospectively untouched Stage~2 test split.
Figure~\ref{fig:method-overview} summarises the intended protocol.

\begin{figure*}[t]
  \centering
  \includegraphics[width=0.98\textwidth]{method_overview.pdf}
  \caption{End-to-end project workflow.  Traditional and deep models share
  the same frozen class mapping and partitions.  Validation controls
  model/configuration selection and train-only pilots control feasibility;
  current artifacts are hash-bound before any future test access.}
  \label{fig:method-overview}
\end{figure*}

Our contributions are fourfold.  First, we construct a deterministic,
audited 500-species protocol with 40 training, 10 validation, and 10 test
images per species.  Second, we provide complete traditional and deep
baselines under common metrics, including training and inference timing.
Third, we compare feature components, ImageNet pretraining, augmentation
strength, and backbone freezing on validation data; the matched ResNet18
augmentation pair provides the strict one-factor ablation.
Fourth, we complement clean evaluation with corruption robustness, confusion
analysis, failure cases, and target-specific Grad-CAM visualisations.  The
result is not merely a leaderboard: it is evidence about which design choices
matter, where performance breaks down, and which apparent explanations remain
uncertain.
